\begin{quote}
\itshape
Governing question: For which buyer, confronting which real alternative, under which beliefs and constraints, is our offering worth a particular sacrifice: and how could we find out without confusing a compliment with a purchase?
\end{quote}

\section*{How This Chapter Fits the Book}

Chapter 14 examined the cost of reliably producing a result. This chapter turns the telescope around. It asks what that result is worth \emph{to the person making the purchase}. Chapter 16 will translate this understanding into pricing architecture, tiers, negotiation, and experiments. Chapter 17 will put the transactions into a complete business model and accounting system. These are connected disciplines, but a founder becomes dangerously confused when they are treated as interchangeable.

This is not a chapter about persuading people to pay more. It is about learning why some exchanges are worth making, why others are rejected, and which unmet jobs can support a business at all. It incorporates and develops the customer-value, psychology, alternatives, and Bayes Institute insights from the earlier essay \emph{The Theory of Pricing and Value Capture}; that earlier essay is a reference, not evidence that any proposed price for Bayes has been validated.

\paragraph*{Epistemic Convention.} All numerical customers, prices, willingness-to-pay estimates, costs, and conversion rates in this chapter are illustrative unless expressly described as externally sourced. Heuristics and scorecards are original management devices, not validated psychometric scales. Economic models assume specified conditions and do not literally describe every purchase.

\section*{Part I: The Decision Inside the Transaction}

\subsection*{1. The Founder's First Mistake: Confusing Worth with Demand}

Imagine a skilled engineer who would benefit enormously from an expert review of a production AI system. A review might prevent a serious mistake. The engineer might praise the idea, tell colleagues about it, and even agree that it is professionally important. Yet when offered a review for ₹12,000, they decline. They are saving for a move, their manager will review the system internally, and the deadline has passed. This is not necessarily evidence that expert review has little value. It is evidence that \emph{this particular exchange}, under \emph{these conditions}, did not win.

Founders often ask, ``What is the value of this product?'' The more precise question is, ``What will this person exchange to obtain this outcome rather than follow the best available alternative?'' Those are radically different questions. The 1st invites stories about intrinsic importance; the 2nd forces us to study budgets, beliefs, substitutes, risk, timing, and actual behavior. A life-changing product can be commercially nonviable if its beneficiaries cannot pay, cannot discover it, cannot trust it, or cannot use it.

A school may confer significant social benefit while collecting insufficient tuition to support itself. An elegant research instrument may have enormous technical merit but only 5 laboratories able to approve its purchase. A workflow tool might generate impressive modeled savings yet fail because the savings accrue to 1 department and the invoice is charged to another. \textbf{Willingness to pay is not an award for having created something important. It is a property of a decision.}

\subsection*{2. Utility: Why Rational Choice Is Broader Than Financial Return}

Economists use \emph{utility} to represent preferences over outcomes. It is not necessarily happiness, cash value, or an observable unit stored in the customer's head. In ordinary choice modeling, utility functions help describe a ranking: a person prefers 1 feasible option to another given their objectives and constraints. The same customer might choose a cheap phone for a child, a premium 1 for a professional need, and a very different product as a gift. We should not dismiss all 3 choices as inconsistent merely because the devices contain similar chips.

The benefits people seek include functional performance, enjoyment, relief from uncertainty, identity, affiliation, aesthetics, convenience, safety, learning, and future options. The costs include money, foregone alternatives, effort, embarrassment, delay, switching risk, and exposure to failure. Some benefits can be monetized with reasonable confidence. Others cannot, and invented currency conversions create only apparent precision.

A useful conceptual model is:

\begin{equation}
U_{ij} = B_{ij} - S_{ij},
\end{equation}
for person $i$ and option $j$, where $B_{ij}$ represents anticipated multidimensional benefits and $S_{ij}$ total anticipated sacrifice. The subtraction is a reminder of trade-offs, not a claim that all terms share a single cardinal currency. Purchasing requires more than finding a high utility: customers must know the option exists, possess enough money or authority, and trust that the relevant benefits are attainable. The founder should understand \emph{decision utility}, \emph{experienced utility}, and \emph{remembered utility} as potentially different objects. People can buy something confidently and regret it afterward; they can also resist a difficult learning experience that later proves immensely valuable.

\subsection*{3. Willingness to Pay Is a Counterfactual, Not a Sticker Number}

For a precisely defined offer and situation, imagine raising the price incrementally while keeping other conditions constant. The customer's reservation price is the largest price at which they would still prefer the offer to their best feasible alternative. We call this \emph{maximum willingness to pay} (WTP). It is helpful in theory but rarely directly observable. Often the person does not know their own threshold until confronted with a real choice, and choices can be noisy.

WTP also depends on the counterfactual. Suppose a company would pay up to ₹4 lakh for a fraud-review tool if the only alternative were expensive manual examination. If an adequate competing product is available at ₹1 lakh, the incremental WTP for the 1st vendor may collapse. The same functionality remains, but the \emph{economic context} changes. Conversely, a specialized vendor may earn a premium when its reliability, audit trail, or integration changes a high-stakes decision.

The founder needs to separate 4 distinct statements: ``The customer could receive benefit $B$''; ``The customer \emph{expects} benefit $B$''; ``The customer expects \emph{incremental} benefit $B$ over the next-best option''; and ``The customer will make a real purchase at price $P$.'' The distance between these statements is where much entrepreneurial optimism dies. A credible business is built by reducing that distance with evidence.
